% BHU PYQ -- Manually transcribed & error-corrected
% Paper: BCH-312 : Income Tax Law and Accounts
% Exam: B.Com. (Hons.) Semester V Examination, 2022-23

\documentclass[11pt,a4paper]{article}
\usepackage[a4paper,top=2.5cm,bottom=2.5cm,left=2.8cm,right=2.8cm,headheight=14pt]{geometry}
\usepackage{amsmath,amssymb}
\usepackage[T1]{fontenc}
\usepackage[utf8]{inputenc}
\usepackage{lmodern,microtype,fancyhdr,enumitem,hyperref,lastpage,parskip}

\hypersetup{
    colorlinks=true,
    urlcolor=black,
    linkcolor=black,
    pdftitle={BCH-312: Income Tax Law and Accounts -- BHU B.Com. (Hons.) Sem V 2022-23}
}

\pagestyle{fancy}
\fancyhf{}
\renewcommand{\headrulewidth}{0.4pt}
\renewcommand{\footrulewidth}{0.4pt}
\fancyhead[L]{\small   \textit{Banaras Hindu University}}
\fancyhead[R]{\small   \textit{BCH-312: Income Tax Law and Accounts}}
\fancyfoot[C]{\small Page  \thepage\ of \pageref{LastPage}}


\newlist{parts}{enumerate}{2}
\setlist[parts,1]{label=(\alph*),leftmargin=*,itemsep=5pt,topsep=3pt}
\setlist[parts,2]{label=(\roman*),leftmargin=*,itemsep=3pt,topsep=2pt}

\newcommand{\pts}[1]{\hfill   \textit{[#1]}}


\begin{document}


\begin{center}
  {\large\bfseries BANARAS HINDU UNIVERSITY}\\[2pt]
  {
\normalsize B.Com. (Hons.) --- Semester V Examination, 2022-23}\\[6pt]
  \rule{0.8\linewidth}{0.5pt}\\[6pt]
  {\large\bfseries Paper No. BCH-312: Income Tax Law and Accounts}\\[4pt]
  {
\normalsize Time: 3 Hours\hfill Maximum Marks: 70}
\end{center}

\rule{\linewidth}{0.4pt}

\smallskip



\noindent   \textit{Instructions: Attempt all questions. All questions carry equal marks (14 marks each).}

\bigskip




\noindent   \textbf{Question 1.} \\
The following are the incomes of Mr. 'X' for the previous year 2021-22:

\begin{enumerate}[label=(\arabic*),leftmargin=*]
  \item Dividend from Indian Company. \hfill Rs. 10,000
  \item Profit from Business in Sri Lanka received in India. \hfill Rs. 12,000
  \item Profit from Business in Nepal deposited in Bank there. This business is controlled from India. \hfill Rs. 20,000
  \item Profit from business in Varanasi (Controlled by Canada Head Office). \hfill Rs. 11,000
  \item Interest received from a Non-resident Mr. 'Y' on the loan provided to him for a business carried on in India. \hfill Rs. 5,000
  \item Income was earned in Uganda and received there but brought in India. \hfill Rs. 8,000
  \item Share of income from Indian partnership firm. \hfill Rs. 15,000
  \item Rent from house property in India received in America. \hfill Rs. 20,000
  \item Interest on debentures of an Indian Company received in Dubai. \hfill Rs. 3,000
  \item Capital Gain on sale of agricultural land situated in Mirzapur (urban area). \hfill Rs. 8,000
\end{enumerate}
Compute his Gross total income, if he is:

\begin{parts}
  \item Ordinary resident
  \item Not-ordinarily resident
  \item Non-resident
\end{parts}\pts{14}

\centerline{   \textbf{OR}}




\noindent ``Income tax is charged on the income of the previous year''. Do you fully agree with this statement? If not, what are the exceptions?\pts{14}

\medskip\rule{\linewidth}{0.2pt}




\noindent   \textbf{Question 2.} \\
Mr. Shiva who is employed as Manager in a factory in Varanasi and is covered under the Payment of Gratuity Act, 1972 retired on 31 extsuperscript{st} December, 2021 after completing 34 years of service. At the time of retirement the employer paid him a gratuity of Rs. 5,10,000 and the accumulated balance of Recognised Provident Fund Rs. 2,40,000. He is also entitled to a pension of Rs. 7,500 p.m. He commuted his half pension and received on 1 extsuperscript{st} March, 2022 the Commuted Value of pension amounting to Rs. 2,70,000. On the basis of the following information compute his income from salary for the assessment year 2022-23 assuming that salary and pension are due on the last day of the month:

\begin{itemize}[leftmargin=*]
  \item Basic Pay: Rs. 1,89,000 (Rs. 21,000 p.m. for 9 months)
  \item Dearness allowance: Rs. 45,000 (Rs. 5,000 p.m. for 9 months)
  \item Bonus: Rs. 4,000
  \item Project allowance: Rs. 3,000
  \item House rent allowance: Rs. 4,200 p.m.
  \item Rent paid: Rs. 3,000 p.m. for 12 months
  \item Employer's and employee's contribution to Recognized Provident Fund: Rs. 30,000 each
\end{itemize}\pts{14}

\centerline{   \textbf{OR}}


\begin{parts}
  \item How would you determine the annual value of a house property which remained vacant for part of the previous year?\pts{7}
  \item Compute taxable income from house property from the following information (Assessment Year 2022-23):
  
\begin{enumerate}[label=(\roman*),leftmargin=*]
    \item Fair Market Rent \hfill Rs. 80,000
    \item Actual Rent received \hfill Rs. 72,000
    \item Municipal valuation \hfill Rs. 50,000
    \item Standard Rent \hfill Rs. 60,000
    \item Municipal Taxes \hfill 20\%
    \item Interest paid \hfill Rs. 18,000
  \end{enumerate}\pts{7}
\end{parts}

\medskip\rule{\linewidth}{0.2pt}




\noindent   \textbf{Question 3.} \\
Shri 'S' furnishes the following information relevant for the Assessment year 2022-23:

\begin{center}
   \textbf{Profit \& Loss Account for the year ended 31 extsuperscript{st} March, 2022} \\[6pt]
  
\begin{tabular}{|p{4.3cm}|r|p{4.3cm}|r|}
    \hline
   \textbf{Particulars} &   \textbf{Amount (Rs.)} &   \textbf{Particulars} &   \textbf{Amount (Rs.)} \\ \hline
    \raggedright To Office expenses & 45,000 & \raggedright By Gross Profit & 3,43,000 \\
    \raggedright To Sundry expenses & 39,000 & \raggedright By Sundry receipts & 11,000 \\
    \raggedright To Entertainment expenses & 15,000 & \raggedright By Bad debts recovered (Not allowed earlier) & 7,100 \\
    \raggedright To Audit fee & 12,000 & \raggedright By Customs duties recovered from the Government (Allowed earlier as deduction) & 32,500 \\
    \raggedright To Legal charges & 4,000 & \raggedright By Gifts received from father & 1,43,000 \\
    \raggedright To Extension of building & 6,000 & & \\
    \raggedright To Depreciation on Plant \& Machinery & 23,000 & & \\
    \raggedright To Salary to staff & 43,000 & & \\
    \raggedright To Bonus to staff & 36,000 & & \\
    \raggedright To Contribution towards Recognised Provident Fund & 15,000 & & \\
    \raggedright To Contribution towards unapproved Gratuity Fund & 4,000 & & \\
    \raggedright To Provision for Goods and Services Tax & 25,000 & & \\
    \raggedright To Goods and Services Tax & 38,000 & & \\
    \raggedright To Payment to a National laboratory for Scientific research & 49,600 & & \\
    \raggedright To Net Profit & 1,82,000 & & \\ \hline
   \textbf{Total} &   \textbf{5,36,600} &   \textbf{Total} &   \textbf{5,36,600} \\ \hline
  \end{tabular}
\end{center}




\noindent   \textbf{Additional information:}

\begin{parts}
  \item Payment to a National laboratory is for the purpose of carrying on approved scientific research, not related to the business. Besides, 'S' purchases a plant of Rs. 30,000 for the purpose of carrying on Scientific research related to his business. Neither Cost of plant nor depreciation thereon is debited to Profit \& Loss account.
  \item Depreciation on Plant and Machinery and extension of building as per income-tax rule is Rs. 19,000.
  \item Goods and Services Tax of Rs. 38,000 includes interest for late payment of GST Rs. 1,200 and penalty for evading GST Rs. 10,000.
  \item Provision for GST is however paid on July 10, 2022. Evidence of payment is submitted along with the return of income.
  \item Salary to Staff includes a payment of pension of Rs. 8,000 to the widow of a former employee.
\end{parts}
Compute income from business of Shri 'S' for the Assessment year 2022-23.\pts{14}

\centerline{   \textbf{OR}}




\noindent Enumerate expenses which are allowed in computing taxable profits of a business and also state expenses or losses which are not admissible.\pts{14}

\medskip\rule{\linewidth}{0.2pt}




\noindent   \textbf{Question 4.} \\
From the following information compute the taxable capital gains for the assessment year 2022-23:

\begin{parts}
  \item 'A' purchased a house property in May, 2010 for Rs. 16,70,000.
  \item 'A' sold the property on 20 extsuperscript{th} June, 2021 for Rs. 80,00,000.
  \item Stamp Valuation Authority valued the property Rs. 90,00,000 which is not objected by the seller or the buyer of the house.
  \item On 15 extsuperscript{th} November, 2022 'A' purchased a residential flat Rs. 25,00,000.
  \item 'A' deposited Rs. 10,00,000 in Capital Gain Account Scheme, 1988.
\end{parts}\pts{14}

\centerline{   \textbf{OR}}




\noindent Write short notes on the following:

\begin{parts}
  \item Bond Washing Transactions
  \item Family Pension
  \item Deductions U/S 80 D and U/S 80 TTB of Income Tax Act, 1961
\end{parts}\pts{14}

\medskip\rule{\linewidth}{0.2pt}




\noindent   \textbf{Question 5.} \\
From the following information compute Gross Total Tax (old regime) for the Assessment year 2022-23:

\begin{enumerate}[label=(\roman*),leftmargin=*]
  \item Gross Salary \hfill Rs. 7,25,000
  \item Interest on N.S.C. \hfill Rs. 8,000
  \item Interest on Saving Bank Deposit \hfill Rs. 2,000
  \item Expenses on the rehabilitation of handicapped son \hfill Rs. 10,000
  \item Deposited in Recognized Provident Fund \hfill Rs. 15,000
  \item Life Insurance Premium \hfill Rs. 5,000
\end{enumerate}\pts{14}

\centerline{   \textbf{OR}}


\begin{parts}
  \item When does the liability to pay advance tax arise? When such a tax is to be paid?\pts{7}
  \item What is ex-parte assessment? In what circumstances can it be made?\pts{7}
\end{parts}

\vfill
\rule{\linewidth}{0.4pt}

\begin{center}\small
\href{https://bhu-exam.vercel.app/}{Click here to visit our website}\\[4pt]
\href{https://me-aryan.vercel.app/}{Click here to visit our contributor}\\[4pt]
\href{https://quashan-me.vercel.app/}{Click here to visit our contributor}
\end{center}

\end{document}
